\section{Introduction: The Embedded Linux Ecosystem}

One of the defining strengths of the Linux operating system is its modular architecture and unmatched adaptability. When stripped of desktop overhead—such as graphical user interfaces (X11/Wayland), heavy desktop environments, and non-essential background daemons—Linux transforms into a highly compact, resilient operating system capable of running on resource-constrained Single Board Computers (SBCs) and Systems-on-Module (SOMs).

In 2018, an extensive IoT developer survey conducted by the Eclipse Foundation revealed that \textbf{Linux was selected as the primary operating system for IoT and embedded devices by a dominant margin}, eclipsing bare-metal environments, real-time operating systems (RTOS), and proprietary platform OS solutions. 

\begin{conceptbox}{Why Linux Dominates Embedded Systems}
\begin{itemize}[leftmargin=*]
    \item \textbf{High-Level Hardware Abstraction:} Drivers convert raw hardware registers into standard file system nodes (\texttt{/dev}, \texttt{/sys}, \texttt{/proc}), eliminating the need to write bit-manipulation register code for high-level logic.
    \item \textbf{Rich Software Ecosystem:} Access to POSIX APIs, multi-threaded networking stacks (TCP/IP, IPv6, TLS), database engines, web servers, and high-level language runtimes (Python, Node.js, Java, Go).
    \item \textbf{Portability Across Architectures:} Code written against Linux standard C libraries (\texttt{glibc}, \texttt{musl}) compiles seamlessly across ARM, RISC-V, x86, and MIPS hardware platforms.
    \item \textbf{Community \& Vendor Support:} Extensive open-source driver availability for peripheral ICs, Wi-Fi chips, cellular modems, and graphics accelerators.
\end{itemize}
\end{conceptbox}

However, using Linux in embedded hardware requires a fundamental shift in design philosophy. Engineers must choose between competing paradigms: Bare Metal, Real-Time Operating Systems (RTOS), Embedded Linux, or Heterogeneous Multi-Processing (AMP).
